\documentclass[11pt,a4paper]{article}
\usepackage{fontspec,amsmath,amssymb,amsthm,microtype}
\setmainfont{lmroman10-regular.otf}[BoldFont=lmroman10-bold.otf,ItalicFont=lmroman10-italic.otf,BoldItalicFont=lmroman10-bolditalic.otf]
\usepackage{xcolor}
\usepackage[most]{tcolorbox}
\usepackage{mathrsfs}
\providecommand{\lt}{<}
\providecommand{\gt}{>}
\setmonofont{DejaVuSansMono.ttf}[BoldFont=DejaVuSansMono-Bold.ttf,ItalicFont=DejaVuSansMono-Oblique.ttf,BoldItalicFont=DejaVuSansMono-BoldOblique.ttf,Scale=0.86]
\newfontfamily\leanfallback{STIXTwoMath-Regular.otf}[Scale=0.86]
\definecolor{leanrule}{gray}{0.62}
\definecolor{leanback}{gray}{0.962}
\definecolor{leantitle}{gray}{0.22}
\definecolor{leanmuted}{gray}{0.38}
\newtcolorbox{leanbox}[1]{enhanced,breakable,colback=leanback,colframe=leanrule,
  boxrule=0.4pt,arc=0.6mm,left=2.4mm,right=2.4mm,top=1.8mm,bottom=1.8mm,
  before skip=0.8em,after skip=1em,fonttitle=\small\bfseries,coltitle=leantitle,
  colbacktitle=leanback,titlerule=0.3pt,title={#1}}
\newlength{\leancharwidth}
\newenvironment{leancode}{\par\smallskip\begingroup\ttfamily\footnotesize\raggedright
  \settowidth{\leancharwidth}{M}\setlength{\parskip}{0pt}\setlength{\parindent}{0pt}
  \frenchspacing}{\par\endgroup\smallskip}
\newcommand{\leanline}[2]{\hangindent=\dimexpr#1\leancharwidth+4\leancharwidth\relax
  \hangafter=1\noindent\hspace*{#1\leancharwidth}#2\par}
\newcommand{\leansource}[1]{\par{\footnotesize\color{leanmuted}#1}\par}
\newcommand{\leanlabel}[1]{\par\smallskip\noindent{\small\bfseries #1}\par}
\usepackage[margin=25mm]{geometry}
\usepackage[unicode,hidelinks]{hyperref}
\hypersetup{pdfauthor={LangevinLean Project}}
\setlength{\emergencystretch}{3em}
\setlength{\parskip}{0.4em}
\setlength{\parindent}{0pt}
\allowdisplaybreaks[2]
\raggedbottom
\title{\texorpdfstring{The Verlet–Ornstein–Uhlenbeck splitting and its local error}{The Verlet–Ornstein–Uhlenbeck splitting and its local error}}
\author{LangevinLean Project}
\date{}
\begin{document}\maketitle
Draft mathematical exposition; these notes may contain mistakes.


Let \(\mathcal X\) be a real normed space and give \(\mathcal X\times\mathcal X\) the maximum norm \(\|(v,x)\|=\max(\|v\|,\|x\|)\); phase points are written velocity first. Let \(g:\mathcal X\to\mathcal X\) be \(M\)-Lipschitz (in the application \(g=\nabla f\)), \(\gamma\ge0\) the friction and \(h\ge0\) the step. The kinetic drift is \(d_\gamma(v,x)=(\gamma v+g(x),-v)\), so that the kinetic integral equation reads \(Z_t-Z_0=-\int_0^td_\gamma(Z_s)\,ds+(B_t,0)\) for a velocity forcing \(B\), and one explicit Euler step is \(z-hd_\gamma(z)\).




This module defines the Lie splitting step


\[S_h(z,N)=\bigl(e^{-\gamma h}v'+N,\ x'\bigr),\qquad(v',x')=V_h(z),\]


a velocity Verlet step \(V_h\) for the Hamiltonian part followed by exact exponential damping of the velocity and an additive velocity increment \(N\in\mathcal X\). It proves that \(S_h\) is Lipschitz in \(z\) with constant \(1+(M+1)^2h\), that it differs from the kinetic Euler step \(z-hd_\gamma(z)+(N,0)\) by at most \(h^2C(\|z\|+\|g(0)\|)\) with an explicit constant \(C\), and that the filtered increment


\[N_{a,h}=B(a+h)-B(a)-\gamma\int_0^he^{-\gamma(h-s)}\bigl(B(a+s)-B(a)\bigr)\,ds\]


of a forcing path \(B\) differs from the plain increment \(B(a+h)-B(a)\) by at most \(\gamma h\) times the oscillation of \(B\) on \([a,a+h]\). By the identities proved for Brownian segments, \(N_{a,h}\) is exactly the Ornstein–Uhlenbeck velocity response \(\int_a^{a+h}e^{-\gamma(a+h-r)}\,dB_r\) when \(B\) is the Brownian motion. These deterministic estimates are the consistency and stability ingredients of the convergence of the splitting to the kinetic diffusion, proved in the next modules.


\section{\texorpdfstring{Damping the velocity is nonexpansive}{Damping the velocity is nonexpansive}}


\textbf{Definition.} For \(\delta\in\mathbb R\), \(D_\delta(v,x)=(\delta v,x)\).




\textbf{Claim.} If \(0\le\delta\le1\), then \(\|D_\delta z-D_\delta w\|\le\|z-w\|\) for all \(z,w\in\mathcal X\times\mathcal X\).




\textbf{Proof.} Write \(z=(v,x)\), \(w=(v',x')\). Then \(D_\delta z-D_\delta w=(\delta(v-v'),x-x')\), and by the definition of the maximum norm \(\|D_\delta z-D_\delta w\|=\max\bigl(\delta\|v-v'\|,\|x-x'\|\bigr)\), using \(\|\delta u\|=|\delta|\|u\|=\delta\|u\|\). Since \(0\le\delta\le1\), \(\delta\|v-v'\|\le\|v-v'\|\le\|z-w\|\), and \(\|x-x'\|\le\|z-w\|\). A maximum of two numbers bounded by \(\|z-w\|\) is bounded by \(\|z-w\|\). \(\square\)

\begin{leanbox}{Lean: \texttt{velocityDamping}, \texttt{velocityDamping\_\allowbreak{}norm\_\allowbreak{}sub\_\allowbreak{}le}}
\leansource{\texttt{velocityDamping} \quad KineticLangevin/\allowbreak{}KineticSplitting.lean, lines 12–12}
\begin{leancode}
\leanline{0}{def\ velocityDamping\ (δ\ :\ ℝ)\ (z\ :\ E×E)\ :\ E×E\ :=\ (δ\ •\ z.1,z.2)}
\end{leancode}
\leansource{\texttt{velocityDamping\_\allowbreak{}norm\_\allowbreak{}sub\_\allowbreak{}le} \quad KineticLangevin/\allowbreak{}KineticSplitting.lean, lines 16–23}
\begin{leancode}
\leanline{0}{theorem\ velocityDamping\_norm\_sub\_le\ (δ\ :\ ℝ)\ (hδ0\ :\ 0≤δ)\ (hδ1\ :\ δ≤1)}
\leanline{4}{(z\ w\ :\ E×E)\ :\ ‖velocityDamping\ δ\ z-velocityDamping\ δ\ w‖\ ≤\ ‖z-w‖}
\end{leancode}
\leanlabel{How the Lean proof follows the argument}
\texttt{velocity\allowbreak{}Damping \ensuremath{\delta} z} is \texttt{(\ensuremath{\delta} \ensuremath{\bullet} z.\allowbreak{}1,z.\allowbreak{}2)}. In \texttt{velocity\allowbreak{}Damping\_\allowbreak{}norm\_\allowbreak{}sub\_\allowbreak{}le}, \texttt{hv} and \texttt{hx} are the coordinate bounds \texttt{norm\_\allowbreak{}fst\_\allowbreak{}le}, \texttt{norm\_\allowbreak{}snd\_\allowbreak{}le} (each coordinate of a pair has norm at most the norm of the pair); \texttt{change} exposes the maximum norm of the product, \texttt{max\_\allowbreak{}le} splits the goal, and the velocity part uses \texttt{smul\_\allowbreak{}sub}, \texttt{norm\_\allowbreak{}smul} and \texttt{mul\_\allowbreak{}le\_\allowbreak{}of\_\allowbreak{}le\_\allowbreak{}one\_\allowbreak{}left} (\(\delta a\le a\) for \(a\ge0\), \(\delta\le1\)).
\end{leanbox}

\section{\texorpdfstring{The splitting step and its stability}{The splitting step and its stability}}


\textbf{Definition.} With \(V_h=K_{h/2}\circ D'_h\circ K_{h/2}\) the velocity Verlet step (kick \(K_s(v,x)=(v-sg(x),x)\), drift \(D'_h(v,x)=(v,x+hv)\)), let \(S_h(z,N)=\bigl(e^{-\gamma h}v'+N,x'\bigr)\) where \((v',x')=V_h(z)\).




\textbf{Claim.} Let \(g\) be \(M\)-Lipschitz, \(\gamma\ge0\) and \(0\le h\le1\). Then for all \(z,w\in\mathcal X\times\mathcal X\) and \(N\in\mathcal X\),


\[\|S_h(z,N)-S_h(w,N)\|\le\bigl(1+(M+1)^2h\bigr)\|z-w\| .\]


\textbf{Proof.} The increment \(N\) is the same in both steps and cancels: \(S_h(z,N)-S_h(w,N)=D_\delta(V_hz)-D_\delta(V_hw)\) with \(\delta=e^{-\gamma h}\). Since \(\gamma h\ge0\), \(0<\delta\le1\), so by the previous step \(\|D_\delta(V_hz)-D_\delta(V_hw)\|\le\|V_hz-V_hw\|\). The stability of velocity Verlet, proved in the module on Hamiltonian convergence (by composing the Lipschitz bounds \(1+sM\) of the kicks and \(1+h\) of the drift), gives \(\|V_hz-V_hw\|\le(1+(M+1)^2h)\|z-w\|\) for \(0\le h\le1\). \(\square\)

\begin{leanbox}{Lean: \texttt{kineticSplitStep}, \texttt{kineticSplitStep\_\allowbreak{}norm\_\allowbreak{}sub\_\allowbreak{}le}}
\leansource{\texttt{kineticSplitStep} \quad KineticLangevin/\allowbreak{}KineticSplitting.lean, lines 27–28}
\begin{leancode}
\leanline{0}{def\ kineticSplitStep\ (g\ :\ E\ →\ E)\ (γ\ h\ :\ ℝ)\ (z\ :\ E×E)\ (N\ :\ E)\ :\ E×E\ :=}
\leanline{2}{ornsteinStep\ (Real.exp\ (-γ*h))\ 1\ N\ (verletStep\ g\ h\ z)}
\end{leancode}
\leansource{\texttt{kineticSplitStep\_\allowbreak{}norm\_\allowbreak{}sub\_\allowbreak{}le} \quad KineticLangevin/\allowbreak{}KineticSplitting.lean, lines 32–43}
\begin{leancode}
\leanline{0}{theorem\ kineticSplitStep\_norm\_sub\_le\ (g\ :\ E\ →\ E)\ (M\ :\ ℝ≥0)\ (hg\ :\ LipschitzWith\ M\ g)}
\leanline{4}{(γ\ h\ :\ ℝ)\ (hγ\ :\ 0≤γ)\ (hh\ :\ 0≤h)\ (hh1\ :\ h≤1)\ (z\ w\ :\ E×E)\ (N\ :\ E)\ :}
\leanline{4}{‖kineticSplitStep\ g\ γ\ h\ z\ N-kineticSplitStep\ g\ γ\ h\ w\ N‖\ ≤}
\leanline{6}{(1+((M\ :\ ℝ)+1)\textasciicircum{}2*h)*‖z-w‖}
\end{leancode}
\leanlabel{How the Lean proof follows the argument}
\texttt{kinetic\allowbreak{}Split\allowbreak{}Step g \ensuremath{\gamma} h z N} is \texttt{ornstein\allowbreak{}Step (Real.\allowbreak{}exp (-\allowbreak{}\ensuremath{\gamma}*h)) 1 N (verlet\allowbreak{}Step g h z)}: the Ornstein–Uhlenbeck substep with damping \(e^{-\gamma h}\), noise coefficient \(1\) and noise \(N\), applied after \texttt{verlet\allowbreak{}Step}. In \texttt{kinetic\allowbreak{}Split\allowbreak{}Step\_\allowbreak{}norm\_\allowbreak{}sub\_\allowbreak{}le}, \texttt{he} is the cancellation identity, checked componentwise by \texttt{ext <;> simp}; the conclusion chains \texttt{velocity\allowbreak{}Damping\_\allowbreak{}norm\_\allowbreak{}sub\_\allowbreak{}le} (with \(0\le e^{-\gamma h}\le1\) from \texttt{Real.\allowbreak{}exp\_\allowbreak{}pos} and \texttt{Real.\allowbreak{}exp\_\allowbreak{}le\_\allowbreak{}one\_\allowbreak{}iff}) and \texttt{verlet\allowbreak{}Step\_\allowbreak{}norm\_\allowbreak{}sub\_\allowbreak{}le} using \texttt{.\allowbreak{}trans}.
\end{leanbox}

\section{\texorpdfstring{Exponential damping errors without division by the friction}{Exponential damping errors without division by the friction}}


\textbf{Claim.} For \(\gamma\ge0\) and \(h\ge0\),


\[\bigl|e^{-\gamma h}-1+\gamma h\bigr|\le\frac{\gamma^2h^2}2\qquad\text{and}\qquad0\le1-e^{-\gamma h}\le\gamma h .\]


\textbf{Proof.} Recall \(\psi_0(t)=e^{-\gamma t}\), \(\psi_1(t)=\int_0^t\psi_0\), \(\psi_2(t)=\int_0^t\psi_1\), and the identities \(\gamma\psi_1(t)=1-e^{-\gamma t}\) and \(\gamma\psi_2(t)=t-\psi_1(t)\), valid for every real \(\gamma\) and \(t\) (the second is the integral of the first). Multiplying the second by \(\gamma\) and substituting the first,


\[\gamma^2\psi_2(h)=\gamma h-\gamma\psi_1(h)=\gamma h-1+e^{-\gamma h}.\]


For \(h\ge0\), \(\psi_2(h)\ge0\) (integral of a nonnegative function), so the absolute value equals \(\gamma^2\psi_2(h)\). For \(\gamma\ge0\), \(\psi_0\le1\) on \([0,\infty)\), hence \(\psi_1(s)\le s\) and \(\psi_2(h)\le\int_0^hs\,ds=h^2/2\); multiplying by \(\gamma^2\ge0\) gives the first inequality. For the second, \(-\gamma h\le0\) gives \(e^{-\gamma h}\le1\), and the tangent-line inequality \(e^u\ge1+u\) with \(u=-\gamma h\) gives \(1-e^{-\gamma h}\le\gamma h\). \(\square\)




No division by \(\gamma\) occurs, so the estimates include \(\gamma=0\).

\begin{leanbox}{Lean: \texttt{damping\_\allowbreak{}euler\_\allowbreak{}remainder}, \texttt{one\_\allowbreak{}sub\_\allowbreak{}damping\_\allowbreak{}le}}
\leansource{\texttt{damping\_\allowbreak{}euler\_\allowbreak{}remainder} \quad KineticLangevin/\allowbreak{}KineticSplitting.lean, lines 47–57}
\begin{leancode}
\leanline{0}{theorem\ damping\_euler\_remainder\ (γ\ h\ :\ ℝ)\ (hγ\ :\ 0≤γ)\ (hh\ :\ 0≤h)\ :}
\leanline{4}{|Real.exp\ (-γ*h)-1+γ*h|\ ≤\ γ\textasciicircum{}2*h\textasciicircum{}2/2}
\end{leancode}
\leansource{\texttt{one\_\allowbreak{}sub\_\allowbreak{}damping\_\allowbreak{}le} \quad KineticLangevin/\allowbreak{}KineticSplitting.lean, lines 60–64}
\begin{leancode}
\leanline{0}{theorem\ one\_sub\_damping\_le\ (γ\ h\ :\ ℝ)\ (hγ\ :\ 0≤γ)\ (hh\ :\ 0≤h)\ :}
\leanline{4}{0\ ≤\ 1-Real.exp\ (-γ*h)\ ∧\ 1-Real.exp\ (-γ*h)\ ≤\ γ*h}
\end{leancode}
\leanlabel{How the Lean proof follows the argument}
\begin{enumerate}


\item \texttt{damping\_\allowbreak{}euler\_\allowbreak{}remainder}: \texttt{h1} and \texttt{h2} are \texttt{friction\_\allowbreak{}mul\_\allowbreak{}psi\_\allowbreak{}one} and \texttt{friction\_\allowbreak{}mul\_\allowbreak{}psi\_\allowbreak{}two}; \texttt{hid} is the identity \(e^{-\gamma h}-1+\gamma h=\gamma^2\psi_2(h)\), obtained by \texttt{nlinarith} after rewriting with \texttt{psi\_\allowbreak{}zero}; \texttt{hnon} is \texttt{psi\_\allowbreak{}nonneg} and \texttt{hbound} multiplies \texttt{psi\_\allowbreak{}two\_\allowbreak{}le} by \(\gamma^2\). The absolute value is removed with \texttt{abs\_\allowbreak{}of\_\allowbreak{}nonneg}.

\item \texttt{one\_\allowbreak{}sub\_\allowbreak{}damping\_\allowbreak{}le} proves the two components with \texttt{Real.\allowbreak{}exp\_\allowbreak{}le\_\allowbreak{}one\_\allowbreak{}iff} and \texttt{Real.\allowbreak{}add\_\allowbreak{}one\_\allowbreak{}le\_\allowbreak{}exp}.

\end{enumerate}
\end{leanbox}

\section{\texorpdfstring{The damped Hamiltonian Euler stage versus kinetic Euler}{The damped Hamiltonian Euler stage versus kinetic Euler}}


\textbf{Claim.} Let \(g\) be \(M\)-Lipschitz, \(\gamma\ge0\), \(h\ge0\), \(\delta=e^{-\gamma h}\). Then for every \(z\in\mathcal X\times\mathcal X\)


\[\bigl\|D_\delta\bigl(z-hd_0(z)\bigr)-\bigl(z-hd_\gamma(z)\bigr)\bigr\|\le h^2\Bigl(\frac{\gamma^2}2+\gamma(M+1)\Bigr)\bigl(\|z\|+\|g(0)\|\bigr).\]


\textbf{Proof.} Write \(z=(v,x)\). Then \(z-hd_0(z)=(v-hg(x),x+hv)\), whose damped version is \((\delta v-\delta hg(x),x+hv)\), and \(z-hd_\gamma(z)=(v-\gamma hv-hg(x),x+hv)\). The position components agree, so


\[D_\delta\bigl(z-hd_0(z)\bigr)-\bigl(z-hd_\gamma(z)\bigr)=\bigl((\delta-1+\gamma h)v+(1-\delta)h\,g(x),\ 0\bigr),\]


and its norm is the norm of the velocity component. By the Lipschitz property and the triangle inequality, \(\|g(x)\|\le\|g(x)-g(0)\|+\|g(0)\|\le M\|x\|+\|g(0)\|\le M\|z\|+\|g(0)\|\). Using the triangle inequality, the previous step (\(|\delta-1+\gamma h|\le\gamma^2h^2/2\), \(0\le1-\delta\le\gamma h\)) and \(\|v\|\le\|z\|\),


\[\begin{aligned}
&\bigl\|(\delta-1+\gamma h)v+(1-\delta)hg(x)\bigr\|\\
&\quad\le\frac{\gamma^2h^2}2\|z\|
  +\gamma h\cdot h\bigl(M\|z\|+\|g(0)\|\bigr)\\
&\quad=h^2\Bigl(\bigl(\tfrac{\gamma^2}2+\gamma M\bigr)\|z\|
  +\gamma\|g(0)\|\Bigr)\\
&\quad\le h^2\Bigl(\frac{\gamma^2}2+\gamma(M+1)\Bigr)
  \bigl(\|z\|+\|g(0)\|\bigr).
\end{aligned}\]


where the last inequality holds because the difference of the two sides is \(h^2\bigl[(\frac{\gamma^2}2+\gamma M)\|g(0)\|+\gamma\|z\|\bigr]\ge0\). \(\square\)

\begin{leanbox}{Lean: \texttt{damping\_\allowbreak{}hamiltonianEuler\_\allowbreak{}defect}}
\leansource{\texttt{damping\_\allowbreak{}hamiltonianEuler\_\allowbreak{}defect} \quad KineticLangevin/\allowbreak{}KineticSplitting.lean, lines 68–108}
\begin{leancode}
\leanline{0}{theorem\ damping\_hamiltonianEuler\_defect\ (g\ :\ E\ →\ E)\ (M\ :\ ℝ≥0)\ (hg\ :\ LipschitzWith\ M\ g)}
\leanline{4}{(γ\ h\ :\ ℝ)\ (hγ\ :\ 0≤γ)\ (hh\ :\ 0≤h)\ (z\ :\ E×E)\ :}
\leanline{4}{‖velocityDamping\ (Real.exp\ (-γ*h))\ (z-h\ •\ kineticNegativeDrift\ g\ 0\ 1\ z)-}
\leanline{6}{(z-h\ •\ kineticNegativeDrift\ g\ γ\ 1\ z)‖\ ≤}
\leanline{6}{h\textasciicircum{}2*(γ\textasciicircum{}2/2+γ*((M\ :\ ℝ)+1))*(‖z‖+‖g\ 0‖)}
\end{leancode}
\leanlabel{How the Lean proof follows the argument}
In \texttt{damping\_\allowbreak{}hamiltonian\allowbreak{}Euler\_\allowbreak{}defect} the drift is \texttt{kinetic\allowbreak{}Negative\allowbreak{}Drift g \ensuremath{\gamma} 1}, i.e. \(d_\gamma\) with inverse mass one. \texttt{hgrowth} is the linear growth bound of \(g\); \texttt{hid} is the componentwise identity (each component checked by \texttt{module}); \texttt{Prod.\allowbreak{}norm\_\allowbreak{}mk} and \texttt{max\_\allowbreak{}eq\_\allowbreak{}left} reduce the maximum norm of \((u,0)\) to \(\|u\|\). \texttt{he} and \texttt{h\ensuremath{\delta}} are the two estimates of the previous step, \texttt{hs} the triangle inequality, and \texttt{ha}, \texttt{hb} bound its two terms with \texttt{mul\_\allowbreak{}le\_\allowbreak{}mul}. The \texttt{calc} block carries out the displayed chain.
\end{leanbox}

\section{\texorpdfstring{The local defect of the splitting step}{The local defect of the splitting step}}


\textbf{Claim.} Let \(g\) be \(M\)-Lipschitz, \(\gamma\ge0\) and \(0\le h\le1\). Then for all \(z\in\mathcal X\times\mathcal X\) and \(N\in\mathcal X\)


\[\bigl\|S_h(z,N)-\bigl(z-hd_\gamma(z)+(N,0)\bigr)\bigr\|\le h^2C\bigl(\|z\|+\|g(0)\|\bigr),\qquad C=(M+1)(M+2)+\frac{\gamma^2}2+\gamma(M+1).\]


\textbf{Proof.} The increment \(N\) appears in the velocity of both terms and cancels, so the left-hand side equals \(\|D_\delta(V_hz)-(z-hd_\gamma(z))\|\) with \(\delta=e^{-\gamma h}\). Insert the intermediate point \(D_\delta(z-hd_0(z))\) and use the triangle inequality:


\[\bigl\|D_\delta(V_hz)-(z-hd_\gamma z)\bigr\|\le\bigl\|D_\delta(V_hz)-D_\delta(z-hd_0z)\bigr\|+\bigl\|D_\delta(z-hd_0z)-(z-hd_\gamma z)\bigr\|.\]


By the first step and the Verlet defect estimate of the module on Hamiltonian convergence (valid for \(0\le h\le1\)), the first term is at most \(\|V_hz-(z-hd_0z)\|\le h^2(M+1)(M+2)(\|z\|+\|g(0)\|)\). The second term is bounded by the previous step. Adding the two bounds gives the claim. \(\square\)




The increment \(N\) is kept exactly; only the deterministic part of the step is compared with Euler.

\begin{leanbox}{Lean: \texttt{kineticSplitStep\_\allowbreak{}euler\_\allowbreak{}defect}}
\leansource{\texttt{kineticSplitStep\_\allowbreak{}euler\_\allowbreak{}defect} \quad KineticLangevin/\allowbreak{}KineticSplitting.lean, lines 112–129}
\begin{leancode}
\leanline{0}{theorem\ kineticSplitStep\_euler\_defect\ (g\ :\ E\ →\ E)\ (M\ :\ ℝ≥0)\ (hg\ :\ LipschitzWith\ M\ g)}
\leanline{4}{(γ\ h\ :\ ℝ)\ (hγ\ :\ 0≤γ)\ (hh\ :\ 0≤h)\ (hh1\ :\ h≤1)\ (z\ :\ E×E)\ (N\ :\ E)\ :}
\leanline{4}{‖kineticSplitStep\ g\ γ\ h\ z\ N-(z-h\ •\ kineticNegativeDrift\ g\ γ\ 1\ z+(N,0))‖\ ≤}
\leanline{6}{h\textasciicircum{}2*(((M\ :\ ℝ)+1)*((M\ :\ ℝ)+2)+γ\textasciicircum{}2/2+γ*((M\ :\ ℝ)+1))*(‖z‖+‖g\ 0‖)}
\end{leancode}
\leanlabel{How the Lean proof follows the argument}
In \texttt{kinetic\allowbreak{}Split\allowbreak{}Step\_\allowbreak{}euler\_\allowbreak{}defect}, \texttt{he} is the cancellation of \(N\); \texttt{htri} is the triangle inequality \texttt{norm\_\allowbreak{}sub\_\allowbreak{}le\_\allowbreak{}norm\_\allowbreak{}sub\_\allowbreak{}add\_\allowbreak{}norm\_\allowbreak{}sub} through the intermediate point; \texttt{h1} chains \texttt{velocity\allowbreak{}Damping\_\allowbreak{}norm\_\allowbreak{}sub\_\allowbreak{}le} with \texttt{verlet\allowbreak{}Step\_\allowbreak{}euler\_\allowbreak{}defect}; \texttt{h2} is \texttt{damping\_\allowbreak{}hamiltonian\allowbreak{}Euler\_\allowbreak{}defect}. The sum of the two bounds is identified with the stated right-hand side by \texttt{ring}.
\end{leanbox}

\section{\texorpdfstring{The filtered increment of a forcing path}{The filtered increment of a forcing path}}


\textbf{Definition.} For a path \(B:\mathbb R\to\mathcal X\), a starting time \(a\) and a step \(h\),


\[N_{a,h}=B(a+h)-B(a)-\gamma\int_0^he^{-\gamma(h-s)}\bigl(B(a+s)-B(a)\bigr)\,ds .\]


(As always for the Bochner integral, the integral of a function that is not integrable is taken to be \(0\); in the applications \(B\) is continuous.) When \(B\) is the Brownian motion, the integration-by-parts identity of the module on Brownian segment equations shows that \(N_{a,h}\) equals the Ornstein–Uhlenbeck response \(\int_a^{a+h}e^{-\gamma(a+h-r)}\,dB_r\).




\textbf{Claim.} Let \(\gamma\ge0\), \(h\ge0\), \(\varepsilon\ge0\), and suppose \(\|B(a+s)-B(a)\|\le\varepsilon\) for all \(s\in[0,h]\). Then


\[\bigl\|N_{a,h}-\bigl(B(a+h)-B(a)\bigr)\bigr\|\le\gamma h\varepsilon .\]


\textbf{Proof.} By definition \(N_{a,h}-(B(a+h)-B(a))=-\gamma\int_0^he^{-\gamma(h-s)}(B(a+s)-B(a))\,ds\), whose norm is \(\gamma\) times the norm of the integral, since \(\gamma\ge0\). For \(s\in(0,h]\) we have \(-\gamma(h-s)\le0\), hence \(0<e^{-\gamma(h-s)}\le1\), and the integrand has norm at most \(\|B(a+s)-B(a)\|\le\varepsilon\). The norm of an integral over an interval of length \(h\) of a function bounded in norm by \(\varepsilon\) is at most \(\varepsilon h\) (this holds also when the integral is \(0\) by convention). Therefore the left-hand side is at most \(\gamma\varepsilon h\). \(\square\)




In the convergence proof \(\varepsilon\) is the modulus of continuity of \(B\) at scale \(h\), so the difference between the filtered and the plain increments is \(o(h)\) uniformly in \(a\) on a bounded horizon.

\begin{leanbox}{Lean: \texttt{ouInputIncrement}, \texttt{ouInputIncrement\_\allowbreak{}defect}}
\leansource{\texttt{ouInputIncrement} \quad KineticLangevin/\allowbreak{}KineticSplitting.lean, lines 133–134}
\begin{leancode}
\leanline{0}{def\ ouInputIncrement\ (γ\ :\ ℝ)\ (B\ :\ ℝ\ →\ E)\ (a\ h\ :\ ℝ)\ :\ E\ :=}
\leanline{2}{B\ (a+h)-B\ a\ -\ γ\ •\ ∫\ s\ in\ 0..h,\ Real.exp\ (-γ*(h-s))\ •\ (B\ (a+s)-B\ a)}
\end{leancode}
\leansource{\texttt{ouInputIncrement\_\allowbreak{}defect} \quad KineticLangevin/\allowbreak{}KineticSplitting.lean, lines 138–156}
\begin{leancode}
\leanline{0}{theorem\ ouInputIncrement\_defect\ (γ\ :\ ℝ)\ (hγ\ :\ 0≤γ)\ (B\ :\ ℝ\ →\ E)}
\leanline{4}{(a\ h\ ε\ :\ ℝ)\ (hh\ :\ 0≤h)\ (hε\ :\ 0≤ε)}
\leanline{4}{(hmod\ :\ ∀\ s\ ∈\ Icc\ 0\ h,\ ‖B\ (a+s)-B\ a‖\ ≤\ ε)\ :}
\leanline{4}{‖ouInputIncrement\ γ\ B\ a\ h-(B\ (a+h)-B\ a)‖\ ≤\ γ*h*ε}
\end{leancode}
\leanlabel{How the Lean proof follows the argument}
\texttt{ou\allowbreak{}Input\allowbreak{}Increment \ensuremath{\gamma} B a h} is the displayed expression, written with the interval integral \texttt{\ensuremath{\int} s in 0.\allowbreak{}.\allowbreak{}h, …}. In \texttt{ou\allowbreak{}Input\allowbreak{}Increment\_\allowbreak{}defect} the hypothesis \texttt{hmod} is the oscillation bound on \texttt{Icc 0 h}. \texttt{hid} isolates the integral term (by \texttt{module}); \texttt{norm\_\allowbreak{}neg}, \texttt{norm\_\allowbreak{}smul} and \texttt{abs\_\allowbreak{}of\_\allowbreak{}nonneg h\ensuremath{\gamma}} extract the factor \(\gamma\); \texttt{interval\allowbreak{}Integral.\allowbreak{}norm\_\allowbreak{}integral\_\allowbreak{}le\_\allowbreak{}of\_\allowbreak{}norm\_\allowbreak{}le\_\allowbreak{}const} gives the bound \(\varepsilon\,|h-0|\), its pointwise hypothesis being proved on \texttt{u\allowbreak{}Ioc 0 h} \(=(0,h]\) (after \texttt{u\allowbreak{}Ioc\_\allowbreak{}of\_\allowbreak{}le}) from \texttt{Real.\allowbreak{}exp\_\allowbreak{}le\_\allowbreak{}one\_\allowbreak{}iff} and \texttt{hmod}; \texttt{nlinarith} concludes.
\end{leanbox}

\end{document}
